\exercisesection{\S~4}

The Lie algebras considered in this paragraph are not necessarily finite dimensional.

\begin{enumerate}
\def\labelenumi{\arabic{enumi})}
\tightlist
\item
  Retain the notations of no. 2. Let \(\lambda \in k^{\mathrm{B}}\) . Associate to any \(\alpha \in \mathrm{B}\) the endomorphisms \(X_{-\alpha}^{\lambda}, H_{\alpha}^{\lambda}, X_{\alpha}^{\lambda}\) of the vector space E such that
\end{enumerate}

\[
X _ {- \alpha} ^ {\lambda} (\alpha_ {1}, \ldots , \alpha_ {n}) = (\alpha , \alpha_ {1}, \ldots , \alpha_ {n})
\]

\[
H _ {\alpha} ^ {\lambda} (\alpha_ {1}, \ldots , \alpha_ {n}) = (\lambda (\alpha) - \sum_ {i = 1} ^ {n} n (\alpha_ {i}, \alpha)) (\alpha_ {1}, \ldots , \alpha_ {n}).
\]

The vector \(X_{\alpha}^{\lambda}(\alpha_{1},\ldots,\alpha_{n})\) is defined by induction on n by the formula

\[
X _ {\alpha} ^ {\lambda} (\alpha_ {1}, \ldots , \alpha_ {n}) = (X _ {- \alpha_ {1}} ^ {\alpha} X _ {\alpha} ^ {\lambda} - \delta_ {\alpha , \alpha_ {1}} H _ {\alpha} ^ {\lambda}) (\alpha_ {2}, \ldots , \alpha_ {n}),
\]

where \(\delta_{\alpha, \alpha_1}\) is the Kronecker symbol; we agree that, if \((\alpha_1, \ldots, \alpha_n)\) is the empty word, then \(X_{\alpha}^{\lambda}(\alpha_1, \ldots, \alpha_n)\) is zero.

Show that Lemmas 1 and 2 remain true for these endomorphisms. We thus obtain a representation \(\rho_{\lambda}:\mathfrak{a}\to \mathfrak{gl}(\mathrm{E})\) such that

\[
\rho_ {\lambda} (x _ {\alpha}) = X _ {\alpha} ^ {\lambda}, \rho_ {\lambda} (h _ {\alpha}) = H _ {\alpha} ^ {\lambda}, \rho_ {\lambda} (x _ {- \alpha}) = X _ {- \alpha} ^ {\lambda}.
\]

¶ 2) Retain the notations of no. 3. Let m be an ideal of a.

\begin{enumerate}
\def\labelenumi{\alph{enumi})}
\item
  Let \(\alpha\) and \(\beta\) be two distinct elements of B. Assume that (ad \(x_{\alpha})^{\mathrm{N}}x_{\beta}\) belongs to \(\mathfrak{m}\) for N sufficiently large. Show that \(x_{\alpha \beta} \in \mathfrak{m}\) . (Apply the results of §1, no. 2 to \(\mathfrak{a} / \mathfrak{m}\) , provided with a suitable \(\mathfrak{sl}(2,k)\) -module structure.)
\item
  Show that \(n + \theta n\) is the smallest finite codimensional ideal of a. Show that this is also the smallest ideal containing the \((\operatorname{ad} x_{\alpha})^{4}x_{\beta}\) and the \((\operatorname{ad} x_{-\alpha})^{4}x_{-\beta}\) .
\end{enumerate}

\begin{enumerate}
\def\labelenumi{\arabic{enumi})}
\setcounter{enumi}{2}
\tightlist
\item
  Let \((\mathfrak{g},\mathfrak{h})\) be a split semi-simple Lie algebra, R the corresponding root system, and B a basis of R. For any \(\alpha \in \mathrm{B}\) (resp. for any pair \((\alpha ,\beta)\in \mathrm{B}^2\) ), let \(\mathrm{R}(\alpha)\) (resp. \(\mathrm{R}(\alpha ,\beta)\) ) be the closed subset of R formed by the \(\pm \alpha\) (resp. the smallest closed subset of R containing \(\pm \alpha\) and \(\pm \beta\) ). Let \(\mathfrak{g}(\alpha)\) (resp. \(\mathfrak{g}(\alpha ,\beta)\) ) be the derived algebra of the algebra \(\mathfrak{h} + \mathfrak{g}^{\mathrm{R}(\alpha)}\) (resp. \(\mathfrak{h} + \mathfrak{g}^{\mathrm{R}(\alpha ,\beta)}\) ), cf.~§3.
\end{enumerate}

\begin{enumerate}
\def\labelenumi{\alph{enumi})}
\item
  Show that \(\mathfrak{g}(\alpha) = kH_{\alpha}\oplus \mathfrak{g}^{\alpha}\oplus \mathfrak{g}^{-\alpha}\) ; it is isomorphic to \(\mathfrak{sl}(2,k)\) .
\item
  Show that \(\mathfrak{g}(\alpha,\beta)\) is semi-simple, and that it is generated by \(\mathfrak{g}(\alpha)\) and \(\mathfrak{g}(\beta)\) . Its root system can be identified with \(\mathrm{R}(\alpha,\beta)\) .
\item
  Let \(\mathfrak{s}\) be a Lie algebra (not necessarily finite dimensional). For all \(\alpha \in \mathrm{B}\) , let \(f_{\alpha}\) be a homomorphism from \(\mathfrak{g}(\alpha)\) to \(\mathfrak{s}\) . Assume that, for any pair \((\alpha, \beta)\) , there exists a homomorphism \(f_{\alpha \beta}: \mathfrak{g}(\alpha, \beta) \to \mathfrak{s}\) that extends both \(f_{\alpha}\) and \(f_{\beta}\) . Show that, in that case, there exists a unique homomorphism \(f: \mathfrak{g} \to \mathfrak{s}\) that extends the \(f_{\alpha}\) . (Use Prop. 4 (i).)
\end{enumerate}

\begin{enumerate}
\def\labelenumi{\arabic{enumi})}
\setcounter{enumi}{3}
\item
  Let \(\mathfrak{g}\) be a splittable semi-simple Lie algebra, and \(\sigma\) an automorphism of \(k\) . Let \(\mathfrak{g}_{\sigma}\) be the Lie algebra obtained from \(\mathfrak{g}\) by extending scalars by means of \(\sigma\) . Show that \(\mathfrak{g}_{\sigma}\) is isomorphic to \(\mathfrak{g}\) . (Use the Cor. of Prop. 4.)
\item
  \begin{enumerate}
  \def\labelenumii{\alph{enumii})}
  \tightlist
  \item
    Let g be a simple Lie algebra, and \(k_{1}\) the commutant of the adjoint representation of g. Show that \(k_{1}\) is a commutative field, a finite extension of k, and that g is an absolutely simple \(k_{1}\) -Lie algebra.
  \end{enumerate}
\end{enumerate}

Conversely, if \(k_{1}\) is a finite extension of k, and g an absolutely simple \(k_{1}\) -Lie algebra, then g is a simple k-Lie algebra, and the commutant of its adjoint representation can be identified with \(k_{1}\) .

\begin{enumerate}
\def\labelenumi{\alph{enumi})}
\setcounter{enumi}{1}
\tightlist
\item
  Let \(k'\) be a Galois extension of k containing \(k_{1}\) . Show that \(\mathfrak{g}_{(k')}\) is a product of \([k_{1}:k]\) absolutely simple algebras. When \(\mathfrak{g}_{(k')}\) is split, these algebras are mutually isomorphic. (Use Exerc. 4.)
\end{enumerate}

\begin{enumerate}
\def\labelenumi{\arabic{enumi})}
\setcounter{enumi}{5}
\tightlist
\item
  Let \(\mathbf{A}\) be a commutative ring, and let \(\mathfrak{u}\) be the A-Lie algebra defined by the family of generators \(\{x, y\}\) and the relations
\end{enumerate}

\[
[ x, [ x, y ] ] = 0, \quad [ y, [ y, [ y, x ] ] ] = 0.
\]

Show that \(\mathfrak{u}\) is a free A-module with basis

\[
\{x, y, [ x, y ], [ y, [ x, y ] ] \}.
\]

Show that, when \(\mathrm{A} = k\) , \(\mathfrak{u}\) is isomorphic to the algebra \(\mathfrak{a}_{+} / \mathfrak{n}\) corresponding to a root system of type \(\mathrm{B}_2\) .

¶7) Let A be a commutative ring in which 2 is invertible, and let u be the A-Lie algebra defined by the family of generators \(\{x,y\}\) and the relations

\[
[ x, [ x, y ] ] = 0, \quad [ y, [ y, [ y, [ y, x ] ] ] ] = 0.
\]

Show that \(\mathfrak{u}\) is a free A-module with basis

\[
\left. \right.\left\{x, y, [ x, y ], [ y, [ x, y ] ], [ y, [ y, [ x, y ] ] ], [ x, [ y, [ y, [ x, y ] ] ] ] \right\}.
\]

Show that, when A = k, u is isomorphic to the algebra \(a_{+}/n\) corresponding to a root system of type \(G_{2}\) .

